\section{Quantum-Secure Trading}
\label{sec:quantum-secure}

Every trade on \Zoo{} \DEX{} settles under a \Quasar{} 4.0 triple-cert
inherited from the \Lux{} L1 \cite{lp-020}. The cert composes three
independent signature layers; a trade is final only when all three
have signed.

\subsection{The triple cert at the trade level}

A finalised trade carries:

\begin{enumerate}[leftmargin=1.4em]
  \item \textbf{$\BLS{}_{12\text{-}381}$ aggregate.} 48 bytes.
        Classical fast-path proof from the \Lux{} validator quorum.
        Useful for the low-latency check that gates pre-final block
        propagation.
  \item \textbf{\Corona{} threshold.} Post-quantum threshold
        proof, Ring-LWE-based. $\sim$2.0--2.4 KB.
  \item \textbf{$\MLDSA{}\text{-}65$ identity proof.} FIPS 204
        conforming, $\sim$3.3 KB per validator. Establishes which
        validator authorised the cert layer.
\end{enumerate}

A trade with only $\BLS{}$ is not final under \Zoo{} \DEX{}; the
matching engine waits for the full triple-cert before marking the
trade as settled.

\subsection{Q-day resilience}

Under a near-term Q-day adversary --- a cryptographically relevant
quantum computer (CRQC) capable of breaking ECDSA and pairing-based
schemes --- the trade record continues to validate because the
\Corona{} and $\MLDSA{}$ proofs do not depend on classical hardness
assumptions. A trade settled on \Zoo{} \DEX{} on April 20, 2026
remains independently verifiable post-Q-day; the chain does not
need a flag day to invalidate pre-Q-day records.

This matters for tokenised securities specifically because
beneficial-ownership records persist for the lifetime of the asset
(potentially decades for fixed-income or equity holdings). A
classical-only signature on a 2026 trade in a 30-year corporate
bond would be retroactively forgeable by a CRQC available in, say,
2034. The triple-cert design eliminates that exposure.

\subsection{Harvest-now-decrypt-later resistance}

A passive observer can record every public on-chain transaction
today and decrypt the classical proofs later when a CRQC becomes
available. This is the harvest-now adversary. The triple-cert
discipline means the harvest-now record is not exploitable: even
once $\BLS{}$ falls, the \Corona{} and $\MLDSA{}$ proofs continue
to authenticate the trade. The settlement record is immune to the
harvest-now class.

\subsection{Validator key management}

Validators run dual-key wallets:

\begin{itemize}[leftmargin=1.2em]
  \item Classical key for $\BLS{}$ operations (fast path).
  \item Threshold $\MLDSA{}\text{-}65$ key for identity proofs.
  \item Threshold \Corona{} key for post-quantum threshold
        operations.
\end{itemize}

Key rotation runs on the LP-017 schedule: every 90 days the
validator set runs a re-keying ceremony on M-Chain \cite{lp-134}.
Stake is anchored to the \emph{validator}, not the per-epoch key,
so rotation is non-disruptive to staked-balance accounting.

\subsection{Trade-finality semantics}

A wallet showing ``filled, final'' to the user is showing the
following on-chain state:

\begin{enumerate}[leftmargin=1.4em]
  \item The trade event is in a block sealed by \Zoo{} D-Chain.
  \item The block hash is referenced inside a \Quasar{}Cert.
  \item The \Quasar{}Cert carries valid signatures on all three
        lanes ($\BLS{}$, \Corona{}, $\MLDSA{}$).
  \item The transfer-agent notify has fired and been acknowledged.
\end{enumerate}

If any of these four conditions fail at the time of display, the
wallet shows ``filled, pending finality''. Time-to-final is sub-1.1s
in normal operation.

\subsection{Light-client verification}

A retail wallet does not have to trust a full node to verify
finality. The triple-cert verifier is small enough to run inside
the wallet:

\begin{lstlisting}[language=C++,caption={Light-client cert verifier}]
bool verify_trade_finality(
    TradeEvent ev,
    QuasarCert4 cert,
    ValidatorSet vs
) {
    // Block contains the trade
    if (!merkle_proof_valid(ev, cert.block_root)) return false;
    // Triple-cert valid against the validator set
    if (!bls_verify(cert.bls, cert.block_root, vs))           return false;
    if (!corona_verify(cert.corona, cert.block_root, vs)) return false;
    if (!mldsa_verify(cert.mldsa, cert.block_root, vs))       return false;
    // Liquidity precompile attested the trade
    if (!liquidity_attestation_valid(ev.attestation)) return false;
    return true;
}
\end{lstlisting}

Practical wallet verification budget is well under 100ms per trade
on consumer hardware (the dominant cost is the $\MLDSA{}$ aggregate
verification scaling with the threshold size).

\subsection{Disaster recovery}

\begin{itemize}[leftmargin=1.2em]
  \item \textbf{Single lattice scheme break.} If a published break
        of either \Corona{} or $\MLDSA{}$ surfaces, trading halts
        pending parameter re-issuance. Existing settled trades
        remain valid because they carry the unbroken lane plus the
        classical lane; new trades require new parameters.
  \item \textbf{Both lattice schemes break.} Trading halts; the
        \DAO{} convenes an emergency response. Pre-existing trades
        retain partial validation (classical lane), but new trades
        are paused until a new PQ scheme is integrated.
  \item \textbf{Classical scheme break only.} Trading continues;
        the classical lane is downgraded to advisory; PQ lanes
        continue to authenticate. This is the steady-state under
        Q-day.
\end{itemize}

\subsection{Why this matters for the asset classes the DEX lists}

Equity, fixed income, RWA, and derivatives all carry long-lived
ownership records. A retail user buying tokenised treasuries with
a 10-year term needs to know the record of that purchase will be
verifiable for the term of the bond. The triple-cert is the
minimum bar for that guarantee.

The \Zoo{} \DEX{}'s value proposition --- ``make money like the 1\%
finally'' (\S\ref{sec:retail-ux}) --- is meaningless without
end-to-end Q-day-resilient settlement. A retail user holding a
30-year corporate bond as a tokenised position must be confident
that the position record is valid when the bond matures, regardless
of cryptographic developments in the intervening decades. The
triple-cert is what makes that promise creditable.
